\section{The issue with Christianity}

I had a different topic in mind, but I decided to address this thoughtful comment first:

\begin{quotex}
I think the issue some take with Christianity lies in the third of your criteria. The Christian religion seems to have been the fountainhead of a large number of subversive movements -- though even as I write this I can anticipate your response will be that Christianity was Judaized. Well, fair enough, but wasn't it also paganized in an earlier period? It seems to me that what is of most value in historical Christianity is what is least essentially Christian about it.
\end{quotex}


We need to look at this from the point of view of Tradition and not in a partisan way. The third criterion is this: ``Does the exterior form provide a safe haven within which the elite can accomplish their tasks?'' Even though the post made no reference to Christianity, the comment nevertheless brought it up.

As a word ``Christianity'' encompasses a long period of history and and too many movements, and yes, most of them subversive. See we are not nominalists, we need to specify something with a determinate content. As a side note, it should be pointed out that the subversive movements are nomilinist, either explicitly or implicitly. For the background to this, I recommend \textit{Ideas Have Consequences}\footnote{\url{https://en.wikipedia.org/wiki/Ideas\_Have\_Consequences}}

So as far as I am concerned, we should only be concerned about the ``Religion of the Middle Ages'' (ROTMA). Some, such as Maurras and Comte regarded it as a form of paganism and Evola considered that some ``rectifications'' were made to make it near pagan. Guenon referred to it as Catholicism, but regarded it --- both esoterically and esoterically --- as an extension of Classical Rome, not of Jerusalem. These are the assumptions of the latter two expounders of Tradition:

\begin{enumerate}


\item The Middle Ages represented a civilization of the Traditional type

\item The end began with the Renaissance and Reformation

\item There is no longer a Traditional religion in the West
\end{enumerate}


So the third point is moot; there exists in the West no Traditional religion readily available, whether Christian or pagan. If this point were the only issue in dispute, it could be intelligently discussed. Unfortunately, discussion is replaced with partisanship, with the most absurd and irrelevant arguments made against Christianity. These come straight from the enemies of the West, so you have to wonder if they are really false flag operations. At some point we can address some of them.

It is not advisable to throw away the accomplishments of the west, described by Donoso Cortes\footnote{\url{http://www.gornahoor.net/?p=223}}so quickly. Because of the ROTMA, Europe had spiritual unity. A knight could travel from Portugal to Poland and encounter the same religion and the same code of chivalry. The attempts nowadays to create a unity based on ``whiteness'' or ``genetic similarity'' are ludicrous and doomed to failure.

If we reject ROTMA out of hand, then we will miss out on what was able to establish a continental unity and Traditional society. Guenon regarded Hermeticism as the esoteric core of ROTMA; this is the perspecitve of the Brahmin. We can also point to the chivalric orders to get the perspective of the kshatriya. Everything said deserves much deeper treatment, and that will be forthcoming.


\flrightit{Posted on 2010-02-26 by Cologero}

\begin{center}* * *\end{center}

\begin{footnotesize}
\begin{sffamily}
\texttt{GFJF on 2010-02-26 at 21:21 said:}

I thought your post `Did Geunon Convert?' was a continuation of `Guenon and Tradition', and that is why I brought up Christianity. While it seems that such wasn't your intention, the mistake has proven to be edifying.

In the latter post you said that ``we would expect Traditionalists to find their home in the Catholic church .'' (Now is a good time to recall that Evola never joined the Catholic church.) In this post you appear to have taken that back; here you say ``there exists in the West no Traditional religion readily available, whether Christian or pagan.'' Although I can see sense in both statements, they do appear to contradict eachother. Perhaps you could (a) qualify them; or if they are not in contradiction, (b) show how you think this is so. I don't make any of these comments in an adversarial spirit.

If memory serves, Maurras believed that the Christian idea contained a dangerous potential for egalitarianism. Though this potential was checked for a time by the ROTMA, the Christian message has proven to be singularly volatile. So even though a paganized Christianity fulfilled the third criterion for around 4-5 hundred years, it still seems to me questionable whether a religion whose popular message is egalitarian really fulfils it.

\hfill

\texttt{Cologero on 2010-02-26 at 23:02 said:}

I'll have to get to Maurras another time, though for most of his life he was a pagan with a high regard for the Catholic Church; it will be worth exploring why. Evola's objections to Christianity were more subtle than crude anti-Catholic bashing; again worth exploring as some of it was incisive, other parts somewhat misguided in my opinion. Ultimately, he felt that its spirit was alien to the Aryan spirit; however, we should not conclude that everything that goes by the name of ``paganism'' in our time is much better.

I don't see Christianity as egalitarian at all. There are celestial hierarchies, which is reflected both in the ecclesiastical hierarchy, chivalric orders, and formerly in the arrangement of society itself. Why that changed is a involved discussion.

My point was, even if somewhat unclear, is that a man adopts the religion of the land; it is only modern individualism that gives us a choice of which religion to take on. Since there is no religion in the land with a full metaphysical teaching, we are in a difficult situation. An elite needs to arise, an elite with actual gnosis. What will arise from that, perhaps a hermeticized Catholicism or a new pagansim? But if you regard medieval Catholicism as a continuation of classical Rome --- as Guenon and I do --- then we can expect a continuation of what we have, even if it superficially appears as a clean break.

``If ... ... truth were recognized by a small number it would still be a result of considerable importance, for it is only in this way that a change of orientation leading to a restoration of the normal order can begin; and that restoration ... will necessarily take place sooner or later'' R Guenon

\hfill

\texttt{Comprehensor on 2010-02-26 at 23:47 said:}

Guenon realized that a restoration of Catholicism was not possible because of severe modernist influences. How much more impossible is it today where V2 Catholicism has become totally politicized by Marxist/materialist elements!

\hfill

\texttt{Matt on 2010-02-27 at 04:37 said:}

Cologero, what are your thoughts on Eastern Orthodoxy and its relation to Tradition? Eastern Orthodox was never mentioned that much by Guenon (at least to me), and Evola never seemed to talk about it in his critique of Christianity.

\hfill

\texttt{Cologero on 2010-02-27 at 18:46 said:}

During the time period I am interested in, there was only one church, orthodoxy and catholicism were the same and represented the continuation of the Roman Empire, even though split in two. At least Orthodoxy pays lip service to the concept of theosis: that is, it is not blind faith that is salvific, but rather one must become ``god-like''. Geunon and Evola talk about it, but few bother to understand. It is not the apotheosis of the ego, but rather the transcendence of the human state.

Fundamental is metaphysical understanding and gnosis. This must precede everything, including religious, racial, political, and social considerations.

\hfill

\texttt{Mark on 2010-03-02 at 03:30 said:}

Let me start out with a quote from a site that uses Nietzsche.

``But there is another ``enemy,'' though unconsciously so, and that is Christianity's own ``will to truth.'' The need to grasp truth, verily the morality of truth, in this sense, is ultimately Christianity's gravest danger. If Christianity had been content to remain a schematism it might still exist; but its will to absolute truth has undermined its credibility, especially in relation to its extravagant competition with science.''

\url{http://www4.hmc.edu:8001/humanities/beckman/nietzsche/Genealogy.htm}


Christianity had its history. I think what has happened in it, was bound to happen. It inherently has this tendency towards desacralization and demythologizing, and because of its positivistic nature, it must with all its force make the death and resurrection a fact that did happen. Hence its ``will to truth'' towards this event, will turn against it.

Just look at the end of your posting, where you state Brahmin and Kshatriyah. The Hindu schematism is still strong, and you are using it, and many Traditionalists use it. I could do this for philosophy as well, the system of Transcendental Idealism can be schematized as the ``veil of maya'', and our machines as ``yantras''. This ability does not exist in Christianity anymore, its ``myth'' does not have this power anymore.

What are your thoughts on this?

\hfill

\texttt{Cologero on 2010-03-03 at 08:31 said:}

Nietzsche was looking at the end result of a series of revolutions, where the spiritual authority was overthrown and eventually lost its own self-understanding. Hence, having lost gnosis or Wisdom, it was left with faith alone. This is in contrast to science; this mock battle has been going on now for centuries. We have no interest in it, although it is a ``hot topic''.

I refer you to this chapter from Evola's ``Pagan Imperialism'':

\textit{Those who know and those who believe}\footnote{\url{http://www.gornahoor.net/PaganImperialism/ThoseWhoKnow.htm}}

\hfill

\texttt{Mark on 2010-03-03 at 13:32 said:}

I do believe that the issue of Christianity is important for understanding the present condition as far as the presence of sacred values or lack thereof. I am looking for my copy The Pagan Temptation by Thomas Molnar, because there is a section at the end of the chapter titled, ``The Christian Desacralization'', where he references both Evola and Guenon. The question is can Christianity in this sense Catholicism regain a sense of spiritual vitality by making the sacred real again in the Europe and America, or was there an inherent flaw that needed to develop out and cause its demise?

This flaw could be this, the intense positivism. Christianity must take its myth as a literal fact, and by doing so, it becomes a question of an empirical truth, not a symbol that points to the transcendent. This then creates the need for philosophical elaboration of the doctrine and dogmas of the church, then this philosophical method turns against it, even in so called Christians. Kant had a miracle critique along with Hume.

Hinduism on the other hand, still has the schematism. Their mythopoetic language is so powerful in explaining things. They can even take modern ideas like virtual reality, and explain it as an allegory to Maya. This computer in front of me is not the ``dead machine'' like it is in the West, but it is a Yantra.

The question is this, is there an inherent flaw in Christianity that so wants the death and resurrection of Jesus to be a historical fact, that it loses sight of the transcendent meaning behind a dying and rising savior? It more and more moves to defending the historicity, and then the skeptics attack, it the vicious cycle repeats. The existentialism of someone like Kierkegaard is the end result of this.

Hypothetical Conversation

Person: I am meditating on the story of Jesus in regards to a counter to evolutionary psychology. The materialist position has it that I am a slave to my genes, and nothing can be done. The ``Christ'' in me shows a divine nature in me that connects to the transcendent.

Modern Christian: Stop talking all of this nonsense about deeper spiritual meaning, this will not get you saved, you must believe that he died and rose literally.

Was this inevitable, due to the nature of Christianity?

\hfill

\texttt{Cologero on 2010-03-04 at 06:11 said:}

Some quick points, because these issues will be dealt with in more detail.

As for the Molnar book, I read a \textit{review}\footnote{\url{http://armchair\_academic.homestead.com/pagantemp.html}}; if the reviewer understands Molnar correctly, the book is disappointing. For example, Schumacher was a Catholic himself and correct in his assertion that there is an intellectual power above reason. Also, it is incorrect to assert that detachment and dispassion are distinctly pagan and foreign to Christian spirituality.

The next point is that there is no reason to rule out miracles out of hand despite Hume and Kant. A historical event may nevertheless have a transcendent significance; ``as above, so below'' ... the historical and contingent are manifestations of something higher.

Again, the debate between a religion based on faith and science is of little consequence; the former is based on opinions about the spiritual world, the latter on opinions about the contingent world. There can be no ultimate agreement on opinions. About gnosis or Wisdom, there is always and necessarily agreement, assuming the levels of understanding are the same.

Do you really need Hinduism to explain away the artifacts of contemporary society?

As the Molnar book shows, the caretakers of the Christian religion have lost direct intuition of the spiritual. There is nothing inherently flawed with the religion. The problem, as documented both by Guenon and Evola, is the degeneration of the castes. Low caste individuals have taken over the positions of leadership in the various Churches and the results can be seen. Was that inevitable? Yes, given the nature of cosmic cycles; no, since there remains the question of the Will.

While the consciousness of Western man may be Christian, the unconscious of Western man remains pagan. Catholicism accomodated that. But it was the North Europeans who believed they could eliminate the pagan element from the Christian spirit and the results have been disastrous. For example, by eliminating the pope and priesthood, they opened the door to egalitarianism and the overthrow of legitimate authority.

\hfill

\texttt{Mark on 2010-03-05 at 07:07 said:}

Follow this as a paragraph by paragraph response, ending with the second to last paragraph

There chapter called the Christian Desacralization goes into issues that you would find interesting, so check it out.

The reference to Kant was that of a person who called himself Christian eliminating aspects of the religion like grace and miracles, because he wanted a ``religion within the limits of reason. It was about the historical process of Christianity's ``will to truth'' that I referenced in the Nietzsche quote.

It is of important consequence if it is inherent within Christianity to have this positivism that works against it in a historical process. I don't believe that a Brahmanical or Kshatriyah like force could have prevented it. I think this is where we disagree.

I was lauding Hinduism in that it is still a powerful schematism. It has a langauge that you even use to explain spiritual ideas. Why did it not go through this ``will to truth'' turning against itself as what happened in Christianity

I think this is the central issue. I hope this gets focused on. I think that this was inherent in the doctrine, not just a cause of ``low caste people''. The focus on the facticity of the death and resurrection will always break through any Hermetic covering. I just read a section on polygenism vs monogenism and the Catholic church, and the Catholic Church is so adamant about having monogenism to protect the concept of original sin. I don't see Christianity as plastic to the forces of Tradition that allow it to be completely rectified. That is the issue

Is Christianity plastic enough so that Traditional Forces are able to make it a proper vessel for Tradition, and have it stay that way?

Or

Is there certain doctrines in it, part of what makes it its kerygmatic core, makes what has happened, destined to happen?

This is what I believe to be the central point, so I will end here.

\hfill

\texttt{Cologero on 2010-03-05 at 09:58 said:}

Kant, Nietzsche, Molnar were not among ``those who know'', so their opinions are of little interest. The ``degeneration of castes'' as an historical process is all too obvious, so your disagreement is the result of a misunderstanding. Obviously, the ``Brahmins'', or what was left of them, did not prevent the process of desacralization. However, it is not because of anything inherent in doctrine (I mean that in the Traditional sense), but here is where we can bring in Nietzsche: forces of resenentment were unleashed that emboldend the masses in the lower castes to rebel against their superiors. Don't forget that there were identifiable elements within Europe that formented and financed such movements.

Your line of reasoning is fruitless. We could also say, for example, that there was something inherent in Classical and Germanic paganism that caused it to yield so quickly and thoroughly to Christianity. Or in Hinduism, because it didn't prevent the Aryan invaders of the sub-continent from self-destruction. But to repeat, I claim that paganism was not lost, but remained in the subconscious of the West and its best parts where incorporated into the ROTMA either through theology and philosophy, or the chivalric system. Rather than concentrating on the failures of this Catholicism, you should instead understand and explain how it sustained a Traditional civilization for such an extended period over such an extended space. Every human movement necessarily deteriorates and requires fresh impulses to re-animate it.

As for Hinduism, it is one traditition among many and I don't regard it as superior to anything in the West. I may use the terminology of the Vedanta for convenience, but the West has everything it needs. See, for example, ``The Shape of Ancient Thought'' by McEvilley.

We don't need to defend all the decisions of the Church hierarchy, and from a certain perspective some may seem to be at odds with Traditional or esoteric doctrine. There are reasons for that, some legitimate, others illegitimate because true gnosis was lost at certain points. If you are really interested in the topic, then you need to look at the development of Hermeticism, which is the hidden core (assuming there was one) of esoterism in the West. If we follow the line from Hermes Trismegistus, Pythagoras, Plato, Plotinus, we find a lot of early Christian theologians at the end.

As for polygenism, some recent writings by Christian Hermetecists make the claim for the existence of ``Pre-Adamic'' people who may be biologically human but not spiritually. I refer you to Jung's ``Answer to Job'' and to ``Gnosis'' by Boris Mouravieff.

You object to the historicity of the resurrection, without giving a reason. Isn't the goal of alchemy the Elixer of Life? So we should not and cannot rule out the overcoming of death \emph{a priori}. There is another book worth checking out: ``Covenant of the Heart'' or ``Lazarus, come forth'' by Valentin Tomberg. He relates forgetting, sleeping, and dying as related processes, with the corresponding opposites: remembering, waking up, and living. The corresponding world teachers are Plato (the Sage who related knowing to remembering), Buddha (the awakend one), and Jesus (who overcame death.) This is how a Hermeticist thinks, so if you want to reject something, then reject this.

Nothing is ``destined'' to happen, as all depends on Wisdom and Will, or ignorance and weakness. I refer you to this though of Joseph de Maistre (another Christian Hermeticist):

\begin{quote}\footnotesize
Very few battles are lost physically. Battles are lost almost always morally; the real victor like the real loser is the one who believes he is.
\end{quote}


\hfill

\texttt{Mark on 2010-03-06 at 08:48 said:}

I am dealing with ideas and the consequences that they have, such as your reference to an earlier post that referenced a book called Ideas have Consequences. You seem to think that it is acceptable to argue that the nominalism of Occam had historical consequences, but not the positivism inherent in Christianity as a whole. So what are the ideas that have consequences, and what are the ideas that don't? Could you give me a list?

Again, it is not a fair comparison with Hinduism and Nordic/Slavic paganism, since I am dealing with ideas, and how well those ideas stand, and only those ideas. Hinduism is still vibrant as far as its language that expresses reality. Traditionalism depends on the language of the metaphysical concepts inherent within Hinduism. Also, Nordic and Slavic paganism did not fail because of the ideology failing, but because a sword was placed to the throat of the pagans. So, if Nordic and Slavic paganism failed because of ideas that failed in the historical process, then I would say that it deserves the same critique as Christianity. But we know that was not the reason.

You use the terminology of Hinduism because of its power. If you stopped using terms like Kshatriyah, Brahmin, Kali Yuga, etc., then your rhetoric would lack its power. Why not use language from only Christianity? Don't use Brahmin and Kshatriyah, use Clergy and Knight, and lets see how well you do in your rhetoric.

I am not trying to come off as rude, but it seems as if you are not dealing with my point. Thomas Molnar, who is a devoutly Catholic philosopher, seems to have dealt with the concepts inherent in Christianity that could lead to this situation we have.

\hfill

\texttt{Cologero on 2010-03-08 at 10:17 said:}

The list of ideas is forthcoming.

The point, however, is not that the ideas of Occam had historical consequences; rather to understand how and why. That is because Occams' metaphysical position is defective.

Again, to be clear, the perspective of this blog is not Traditional Catholicism, but rather metaphysics as taught by the the highest represtentatives of the various Traditions, and brought to light by Guenon (the perspective of the Brahmin) and Evola (the perspective of the Kshatriya). I am not a poet, so I am not concerned with the power of words, but rather of their truth. We use Hindu terminology (1) because it is precise, (2) because it is distant from us --- I find that using European or Christian terms seems to instigate unnecessary polemics, and (3) because there is widespread ignorance about the Traditional forms available in the West. Thoughtful writers in the Middle Ages used their own terminology to describe their social structure quite effectively. This needs to become common knowledge; unfortunately, as part of our dispossession, we denigrate the Medieval period ... all sides are guilty of this.

So, to repeat, I do not regard Molnar as a representative of this perspective and I gave two quite specific reasons. I expect specificity in return.

I need to point out that the Traditional teaching is that the empirical world is a reflection of the spiritual --- as above, so below --- an not merely a ``symbol'' for it. Spiritual reality manifests as empirical reality. Now positivism only accepts the empirical as real, and so deals in facts alone, without reference to the transcendent. This is a low caste point of view, not something a Brahmin or Kshatriya would ever accept. This demonstrates the degeneration of castes: men of low caste have risen to important positions in church hierarchy, men who are unable to understand their own Tradition and can see in it nothing but a collection of empirical facts.

I've provided several references (and I can provide more if pressed) to new contemporary writers who see more than ``mere'' facts; this demonstrates there is nothing defective in the ideas, but instead indicates defects in the ability of many to properly understand them. For the lower castes, believing in a literal resurrection is all they are capable of and they should be left in peace. Others see the spiritual reality behind it. Again to be repetitive, esoteric teaching is not a collection of ideas to be debated or believed in. There are ideas of course, but, they must be alive and directly intuited by the knower in an existential way. For that reason, they cannot be grasped from outside the Tradition, nor by anyone not willing to undergo the necessary transformation or intellectual conversion.

Although this is hardly the entire reason for the propagation of Christianity in Europe, I concede there was often enough a certain amount of compulsion. This caused, as I pointed out, the pagan mind to be forced into the unconscious. Catholicism, to a large extent, incorporated pagan ideas and customs (see, e.g., Clement's Stromata, where we read about Druids, Buddha, Hindusim, and Christian gnosis); the Protestants rejected all that, thus pushing paganism even deeper into the unconscious. Now that it is ``liberated'', this paganism is being expressed in some unhealthy ways. The best of the pagan minds were not beer guzzling libertines.

I am being as clear as I can in the confines of an Internet comment. I don't expect anyone to ``believe'' anything, but I do encourage anyone capable of it to try to develop true gnosis. From that perspective, there is not room for these interminable and irrevolvale polemics.


\hfill

\end{sffamily}
\end{footnotesize}

